\documentclass[a4paper,twoside]{article}
\usepackage{graphicx} % Required for inserting images
\usepackage[hidelinks]{hyperref} % geen kaders om linken
\usepackage{parskip}
\usepackage{bm} % For bolt symbols eg \bm{\nabla}
\usepackage{amsmath, amssymb, amsfonts} % For nicer math format
\usepackage{esint} % For bolletje in integrals
\usepackage{mathrsfs} % For curly arrr
\usepackage{subcaption}
\usepackage[english]{babel}
\usepackage{dsfont}
\usepackage[a4paper, margin=0.24cm, landscape]{geometry} 
% Verander 1.5cm naar bijvoorbeeld 1cm voor een extreem vol blad.
\usepackage[T1]{fontenc}
\usepackage[utf8]{inputenc}
\usepackage{lmodern}
\usepackage{multicol}

%Tikz packages
\usepackage{tikz}
\usetikzlibrary{angles,quotes, babel} % for pic (angle labels)
\usetikzlibrary{arrows.meta,decorations.markings,calc, decorations.pathmorphing, decorations.pathreplacing, positioning, shapes, shadings, patterns}
\usepackage{xcolor}
\usepackage{tikz-3dplot}
\usepackage[siunitx]{circuitikz}

\usepackage{etoolbox} % ifthen
\usepackage[outline]{contour} % glow around text
\usepackage{xfp} % higher precision (16 digits?)
\contourlength{1.1pt}



\usepackage{geometry} %size regulating
\geometry{
 a4paper,
 left=20mm,
 right=20mm,
 top=10mm,
 bottom=20mm,
 }

\begin{document}

\begin{center}
\hrule
\vspace{.2cm}
{\bf {\Huge Samenvatting Thermische Fysica}} \\
\vspace{.2cm}
\end{center}
{\bf T. Cools}  \\
{\bf Tweede Bachelor Fysica en Sterrenkunde}\\
{\bf 2025-2026 }  \\
\hrule

\begin{abstract}
    In deze samenvatting zal per hoofdstuk een overzicht worden gegeven van de belangrijkste concepten, definities, formules, afleidingen en bewijzen die in het vak Thermische Fysica aan bod zijn gekomen.
    De samenvatting is bedoeld voor educatieve doeleinden.
\end{abstract}

\input{Thermische-H1.tex}

\end{document}